% Ad Familiares 13.26 (ad-familiares-13-26) --- English translation
% Translated by Claude (Anthropic) from the Latin (Perseus).
% Style guide: STYLE.md.

\ciceroLetterOpener{CICERO SERVIO S.}

\ciceroSection{1} L. Mescinius is bound to me by the tie that he was my quaestor; and this tie, which --- as I have received it from our ancestors --- I have always reckoned to weigh heavily, he has by his own character and humanity made even more rightly so. And thus I am on such terms with him that I am on more familiar terms with no one, and on no terms more gladly. He seemed to be confident that on his account you would gladly do what you honourably could; but he hoped, even so, that a letter of mine too would carry great weight with you. He so judged the matter for himself, and besides, on the strength of our intimacy, he had often heard me say how agreeable our friendship was, and how close.

\ciceroSection{2} I therefore ask of you --- with that earnestness, naturally, with which you can see that I ought to ask on behalf of a man so closely bound and so intimate to me --- that you unravel and dispatch his business in Achaia, which has come to him as heir of his brother M. Mindius (who carried on trade at Elis), with both the law and the authority you possess, and also with your own personal weight and judgement. For we have laid it down to the men to whom we have entrusted that business that in all matters which come into any dispute they should employ you as arbiter and, where it could be done at your convenience, as adjudicator. I beg you, again and again and most earnestly, that you take this on for my honour's sake.

\ciceroSection{3} There is this further request: if you do not think it inconsistent with your standing, you will do me a most welcome service if --- where certain parties prove more obstinate, and unwilling that the business be settled without litigation --- you remit them, since the case is with a senator, to Rome. So that you may do this with the less hesitation, we have procured letters to you from M. Lepidus the consul, not such as enjoin anything upon you (for we do not consider that to be consistent with your standing), but, so to say, in the manner of letters of recommendation.

\ciceroSection{4} I should write to you how well placed this kindness will be where Mescinius is concerned, were I not confident both that you know this and that I am asking it for myself. For I should like you to consider this --- that I am no less anxious over his affair than he is over his own. But while I am eager for him to come into his own as easily as possible, I am no less anxious that he reckon himself to have gained no small thing by this recommendation of mine.
